\section{讨论与局限性}

本章预先回应审稿人可能提出的质疑，对研究设计中的关键假设、方法边界与已知局限进行自我审查。
\subsection{Meta-label 的主观性}

Difficulty 与 \texttt{model\_issue} 依赖标注者主观判断，这是本研究的内在局限。
缓解措施包括统一培训并提供案例示范，通过 per-tag 共识降低个体偏差，并在 Calibration 上报告 Fleiss' Kappa 等一致性统计以量化噪声水平，计算口径见第~\ref{sec:method-tags-consensus}节。
本文的立场是：即使存在测量噪声，这些标签作为风险信号和场景代理仍具备实用价值。

\subsection{Type 4 的可控性边界（无法保证绝对为 0）}

从审稿视角，能将 ``字段缺失/不合规'' 从事后清洗前移到提交时强约束，是采集协议可控性的加分项。本文通过 UI 侧必填与互斥硬校验尽量使 ``空选'' 与 ``互斥冲突'' 在提交时不可发生，并同步记录被拦截次数作为过程证据。

但我们明确接受一个现实结论：Type 4 无法保证绝对为 0。即便 UI 侧硬阻断，系统仍可能出现 NA 或字段缺失（脚本未加载、客户端/服务端版本不一致、导出缺损、记录被污染等）。因此本文仍保留 Type 4/NA 的定义与审计披露，将其作为口径 T 的过程性质量指标，而不是在最终清洗后 ``消失''。
\subsection{样本量限制}
当前可用候选池为 458 张（\texttt{mp3d\_layout}/test）。在 10 人（标注者规模上限约 15--20 人）与紧周期约束下，主方案的阶段划分与规模见第~\ref{sec:exp-data-org}节，核心主分析使用约 260 张唯一图像（Calibration\_manual 100 + Validation 60 + Test 100）。
同理，标注者规模上限约 15--20 人时，细粒度画像分型容易出现空桶与阈值不稳定。为降低稀疏性与事后解释空间，主分析与主报告使用“可靠度主轴 + 离散风险层级”的 4 类工人画像，并将更细粒度的诊断型分型作为未来工作。
\\ 影响体现在：某些低频场景的 $r_{u}^{(s)}$ 仍可能因为样本数不足或置信区间过宽而不可用。低频场景的处理与退化规则见第~\ref{sec:method-scene-reliability}节；本文在结果中透明报告退化比例，并避免用不稳定的场景估计驱动激进路由。
\\ 补充说明：为避免在小样本下引入额外可调超参数而造成软门控调参质疑，主分析采用保守的 LCB 与退化策略，而不采用需额外超参数的收缩平滑或基于相关性的层级回退；相关扩展仅在未来工作讨论。
\\ 透明披露：在 Results 中报告每个场景的样本量和置信区间宽度。
\subsection{离线路由模拟的局限}
方法是在多标注数据上进行模拟分配后取该 worker 真实标注的离线路由评估。
局限在于这不属于真实在线路由，worker 行为可能因任务分配而改变。
辩护理由：离线模拟是路由研究的标准方法，参考 Kara et al. 和 Liu et al. 的工作；避免重新标注的巨大成本；为未来在线实验提供先验。离线模拟回答的是在固定已采集行为数据下策略的可解释收益下界；路由对标注者行为的正向在线效应如激励与适应性改变属于未来工作范畴，不混入主结论。
\subsection{Convex Hull 的局限性}
凸包外包凹形，可能降低对跨门扩张或凹结构错误的敏感度。本文仅将 $\mathrm{IoU}_{\mathrm{edit}}$ 作为改动幅度的辅助证据，其定义与边界已在第~\ref{sec:method-iou-edit}节说明；主结论优先依赖更几何一致的指标，并在反例筛选中采用多指标组合而非仅依赖 IoU。

\subsection{基本假设与让步}\label{sec:disc-assumptions}
受控组织假设：系统部署在受控组织内，如企业内部标注团队，主要面对噪声、平庸、不稳定而非强对抗的恶意攻击者。
\begin{itemize}
    \item 让步1：在极端 OOD 场景下可能出现多数标注者同时犯错而专家意见被淹没的情况。本文在高 $d_t$ 情形更依赖应激模式下的冗余提升与复核机制；加权共识的用法与边界见第~\ref{sec:method-weighted-consensus}节。
	\item 让步2：循环论证风险的规避，通过 PreScreen 的测试样本与专家参考标准获得 $r_{u}^{(0)}$ 的独立锚点，并由此映射得到 $w_u$，同时采用 leave-one-task-out 的估计协议，避免将同一任务的标签既作为推断依据又作为评价结论。
    \item 让步3：本文不在主分析中估计工人 OOD 韧性曲线这类需要在高 $d_{t}$ 区间大量覆盖的数据驱动函数，例如 $r_{u}(d_{t})$ 的衰减斜率，仅将 $d_{t}$ 作为触发器与分层变量；更精细的韧性建模留待更大规模与更长周期的后续研究。

\end{itemize}

\subsection{未来工作：细粒度工人画像分型}
本文在主分析中使用“可靠度主轴 + 离散风险层级”的 4 类工人画像（$R0/R1/R2/R3$）以保证小样本下的稳健性。若在更大规模的标注团队与更长周期下采集到更丰富的过程数据，可以在保持可靠度主轴不变的前提下，把当前风险层级进一步细化为更诊断化的行为类型，用于更精细的审计与培训反馈。

一种可行的细分方案，是在保持当前离散风险层级可追溯规则不变的前提下，再在每一级内部细分 6 类诊断型画像。其前提是：每类达到最小样本量，且 blind-trust、条件性脆弱与工程失败倾向的估计已经足够稳定，避免空桶导致的过拟合解释。
\begin{itemize}
	\item Reliable：高 $\mathrm{LCB}(r_u)$ 且未触发 blind-trust、条件脆弱或工程失败的核心风险。
	\item Normal：$\mathrm{LCB}(r_u)$ 中等，且仅有弱辅助审计风险，不触发主要风险层级升级。
	\item Trust-vulnerable：全局可靠度尚可，但在 semi 误导性初始化上表现出较高 blind-trust 倾向。
	\item Condition-fragile：全局可靠度尚可，但在高风险桶或特定场景上出现显著退化。
	\item Noisy/Sloppy：可靠度偏低并伴随较高工程失败倾向或持续性低质量过程行为。
	\item Severe unusable：在 manual anchors 或 calibration 中已无法满足最低 admission / main-pool 使用条件。
\end{itemize}
该细分的定位是诊断与解释，而非把小样本下的细类分布当作可泛化结论。主分析仍应优先依赖 $\mathrm{LCB}(r_u)$、离散风险层级与保守路由规则，并在高风险任务上采用更严格的候选池与冗余策略。


